\documentclass[11pt]{article}

\usepackage{acl}
\usepackage{times}
\usepackage{latexsym}
\usepackage{microtype}
\usepackage{inconsolata}
\usepackage{graphicx}
\usepackage{amsmath}
\usepackage{booktabs}
\usepackage{url}

\title{Reproducing and Extending OptIForest for Anomaly Detection}

\author{
  Kishore P \\
  Macquarie University \\
  COMP8240 Applications of Data Science \\
  \texttt{kishore.pandidurai@students.mq.edu.au}
}

\begin{document}

\maketitle

\begin{abstract}
This proposal outlines a reproducibility study of OptIForest, an
unsupervised anomaly detection method based on the isolation forest
mechanism. The proposed project will first reproduce selected results
reported by Xiang et al. using the authors' public Python implementation
and original benchmark data. It will then evaluate OptIForest using
additional existing datasets and newly constructed synthetic anomaly
data. The evaluation will consider detection effectiveness using
AUC-ROC and AUC-PR, computational efficiency, and sensitivity to
branching factor, cut threshold, and sampling size. A preliminary
execution on the supplied AD dataset produced results close to those
reported in the original paper, supporting the feasibility of the
project.
\end{abstract}

\section{Introduction}
Anomaly detection is the task of identifying observations or patterns
that differ substantially from the behaviour expected in a dataset
\citep{chandola2009anomaly}. Such observations may represent events of
practical importance, including fraudulent financial transactions,
cybersecurity intrusions, equipment faults, or unusual health
measurements. In many real applications, anomalies are rare and
reliable labels are expensive to obtain. Unsupervised anomaly detection
is therefore valuable because it attempts to identify unusual
observations without requiring labelled examples during model training.

Isolation Forest is an ensemble method that detects anomalies by
repeatedly partitioning sampled observations using isolation trees
\citep{liu2008isolation}. The method is based on the idea that unusual
observations tend to be separated from the majority using fewer
partitions and consequently have shorter tree paths. Although many
isolation-forest methods use binary trees, Xiang et al. argue that a
binary branching factor is not theoretically optimal
\citep{xiang2023optiforest}. Their proposed OptIForest method introduces
a theoretical measure of isolation efficiency and derives an optimal
average branching factor close to Euler's number, \(e \approx 2.718\).
The practical method also incorporates clustering-based learning to
hash with the aim of improving isolation quality while maintaining
computational efficiency.

This project will investigate the reproducibility and generalisability
of OptIForest. First, selected results will be reproduced using the
authors' public implementation and original benchmark data. The method
will then be evaluated on existing datasets not investigated in the
paper and on newly constructed synthetic anomaly data. The study will
compare detection performance using AUC-ROC and AUC-PR, examine
execution time, and investigate sensitivity to branching factor, cut
threshold, and sampling size. The objective is to determine whether the
published results can be reproduced and whether the paper's claims
continue to hold under new data conditions.

\section{Description of the Selected Paper}

\subsection{Anomaly Detection}
Anomaly detection is the task of identifying observations that differ
substantially from the behaviour expected in a dataset
\citep{chandola2009anomaly}. A normal observation follows a common
pattern found among many observations, whereas an anomaly is rare,
unusual or inconsistent with that pattern. Examples include fraudulent
financial transactions, network intrusions, equipment faults and unusual
medical measurements.

Anomaly detection may be supervised when reliable labelled examples are
available. However, anomaly labels are often rare or expensive to obtain.
Unsupervised anomaly-detection methods therefore attempt to assign
anomaly scores without using class labels during model training. A higher
score normally indicates that an observation is considered more unusual.
Ground-truth labels may still be used after training to evaluate how well
the scores rank genuine anomalies above normal observations.

\subsection{Isolation Forest}
Isolation Forest detects anomalies by isolating observations rather
than constructing a detailed model of normal behaviour
\citep{liu2008isolation}. Each isolation tree is built from a random
sample of the data. At every internal node, the original method randomly
selects a feature and a split value, dividing the observations into
smaller groups. This process continues recursively until observations
are isolated or a stopping condition is reached.

The path length of an observation is the number of divisions required
to reach its terminal node. Anomalies are expected to have shorter path
lengths because their feature values are relatively unusual and they
can often be separated from the main data distribution quickly. Normal
observations are located in denser regions and generally require more
partitions before they are isolated. An ensemble of trees is used to
reduce dependence on any single random tree, and the combined path
lengths are converted into anomaly scores.

For example, consider a dataset in which most financial transactions
range from \$10 to \$300 while one transaction is \$15,000. A randomly
selected split on transaction value is likely to separate the \$15,000
observation after only a small number of divisions. Transactions in the
common range may require many additional divisions. The unusually
short path provides evidence that the \$15,000 transaction should
receive a high anomaly ranking. However, whether it is genuinely
fraudulent must ultimately be determined using contextual evidence or
ground-truth labels.


\subsection{The OptIForest Method}
OptIForest investigates how the branching structure of an isolation
tree affects its ability to separate observations
\citep{xiang2023optiforest}. Let \(v\) represent the branching factor,
\(d\) the tree depth, and \(\psi\) the number of leaf nodes in an
idealised perfect tree. The relationship between these quantities is

\[
\psi = v^d.
\]

The paper defines the isolation area as \(\phi = vd\), representing the
combined effect of tree width and depth. Isolation efficiency is then
defined as

\[
\eta = \frac{\psi}{\phi}.
\]

Under a fixed isolation area, the authors formulate an optimisation
problem that maximises isolation efficiency. Their theoretical solution
gives an optimal branching factor of Euler's number,
\(v=e\approx2.718\), rather than the binary factor \(v=2\) commonly used
by isolation forests. A practical tree cannot create a fractional
number of branches. OptIForest therefore uses a distribution of integer
branching factors, principally binary and ternary branches, whose
expected value approaches \(e\).

OptIForest is built on the multi-branch LSHiForest framework. The method
combines locality-sensitive hashing with clustering-based learning to
hash. Similar observations can be assigned to the same branch, while
the clustering component learns information from the data to improve
partition quality. A cut threshold controls how much of each tree is
constructed using clustering-based learning before the remaining
partitioning is completed using the underlying hashing procedure. The
authors argue that this controlled amount of learning reduces bias
without introducing excessive variance or computational cost.

\subsection{Main Contributions and Claims}
The paper makes three principal contributions
\citep{xiang2023optiforest}. First, it introduces the concept of
isolation efficiency and provides a theoretical analysis suggesting
that the optimal average branching factor is \(e\). This offers an
explanation for why multi-branch isolation structures may outperform
strictly binary trees.

Second, the authors develop OptIForest as a practical unsupervised
anomaly detector. The implementation combines the proposed branching
structure with clustering-based learning to hash and is designed to
balance detection quality with computational efficiency.

Third, the paper evaluates OptIForest through comparative and ablation
experiments. It compares OptIForest with six shallow and deep anomaly
detectors across 20 benchmark datasets. The paper reports the highest
average AUC-ROC and AUC-PR among the compared methods, although
OptIForest is not the best method on every individual dataset. The
ablation studies further report that branching factors close to \(e\)
generally perform well and that the effects of cut threshold and
sampling size vary across datasets. These claims provide specific
results that can be investigated through reproduction and evaluation
on new data.

\section{Evaluation Methodology}

\subsection{Baseline Methods}
OptIForest is compared with six unsupervised anomaly detection
baselines \citep{xiang2023optiforest}. These methods represent
traditional isolation-based approaches, statistical methods, and
deep representation-learning approaches.

\textbf{iForest} is the original Isolation Forest method. It constructs
binary isolation trees using random features and split values
\citep{liu2008isolation}. It provides an important baseline for
determining whether OptIForest's alternative branching structure and
data-dependent learning improve performance.

\textbf{LSHiForest} extends the isolation-forest mechanism using
locality-sensitive hashing and permits multi-branch isolation trees
\citep{zhang2017lshiforest}. It is especially relevant because
OptIForest uses LSHiForest as its underlying framework.

\textbf{ECOD} is a statistical detector that estimates outlier scores
using empirical cumulative distribution functions
\citep{li2022ecod}. It does not require iterative model training and is
deterministic, which explains the zero standard deviations reported for
ECOD in the OptIForest results.

\textbf{REPEN} learns lower-dimensional representations that support
random distance-based outlier detection
\citep{pang2018repen}. It represents a deep learning approach intended
to improve detection in high-dimensional data.

\textbf{RDP} learns representations by predicting randomly generated
distances without requiring anomaly labels \citep{wang2021rdp}. It is
included as another deep representation-learning baseline.

\textbf{DiForest}, or Deep Isolation Forest, combines random neural
representations with the isolation-forest mechanism
\citep{xu2022diforest}. Comparing OptIForest with REPEN, RDP and
DiForest allows the paper to investigate whether its shallower method
can provide competitive detection performance with lower computational
cost.

\subsection{AUC-ROC}
The receiver operating characteristic curve evaluates how well anomaly
scores separate anomalous and normal observations across all possible
decision thresholds \citep{fawcett2006roc}. At each threshold, the
true-positive rate and false-positive rate are calculated as

\[
\mathrm{TPR}=\frac{\mathrm{TP}}{\mathrm{TP}+\mathrm{FN}},
\qquad
\mathrm{FPR}=\frac{\mathrm{FP}}{\mathrm{FP}+\mathrm{TN}},
\]

where TP and TN represent correctly identified anomalous and normal
observations, while FP and FN represent the corresponding errors. The
ROC curve plots TPR against FPR.

AUC-ROC is the area under this curve. A value of 1 represents perfect
ranking, while a value near 0.5 indicates performance similar to random
ranking. Higher values are therefore preferable. A useful
interpretation is the probability that a randomly selected anomaly
receives a higher anomaly score than a randomly selected normal
observation. The OptIForest paper reports the mean and standard
deviation of AUC-ROC across 15 runs.

\subsection{AUC-PR}


The precision--recall curve evaluates anomaly detection performance by
comparing precision and recall across different decision thresholds
\citep{saito2015precision}. Precision measures the proportion of
observations predicted as anomalies that are genuine anomalies, while
recall measures the proportion of all genuine anomalies that are
successfully detected:

\[
\mathrm{Precision} =
\frac{\mathrm{TP}}{\mathrm{TP}+\mathrm{FP}},
\qquad
\mathrm{Recall} =
\frac{\mathrm{TP}}{\mathrm{TP}+\mathrm{FN}}.
\]

AUC-PR summarises the precision--recall relationship across the possible
thresholds. Higher values indicate better anomaly detection performance.
It is particularly useful when anomalies are much less common than normal
observations because it focuses on performance for the rare positive
class. The supplied OptIForest implementation calculates the PR score
using Scikit-learn's average precision function. The paper reports the
mean and standard deviation of this score over 15 experimental runs.

\subsection{Worked Evaluation Example}

Consider five observations ranked from the highest to the lowest anomaly
score. Their ground-truth labels in ranking order are
\([1,0,1,0,0]\), where 1 represents an anomaly and 0 represents a normal
observation. The corresponding anomaly scores are
\([0.95,0.80,0.70,0.40,0.10]\).

There are two anomalies and three normal observations, producing six
possible anomaly--normal pairs. The first anomaly is ranked above all
three normal observations. The second anomaly is ranked above two normal
observations but below one. Therefore, five of the six pairs are ordered
correctly:

\[
\mathrm{AUC\mbox{-}ROC} = \frac{5}{6} \approx 0.833.
\]

The implementation reports average precision as its PR score. The
anomalies occur at ranks 1 and 3. Precision at these ranks is \(1/1\)
and \(2/3\), respectively. Therefore,

\[
\mathrm{Average\ Precision}
= \frac{1}{2}\left(1+\frac{2}{3}\right)
\approx 0.833.
\]

The example shows that both measures reward placing anomalies above
normal observations, while average precision concentrates on the quality
of the rankings at positions containing anomalies.

\subsection{Repeated Experiments and Execution Time}
The original study repeats each experiment 15 times and reports the mean
and standard deviation of its evaluation scores
\citep{xiang2023optiforest}. Repetition is important because random
sampling and randomised tree construction can produce different results.
The mean represents typical performance, while a smaller standard
deviation indicates more stable performance. Training and testing times
are also recorded to evaluate computational efficiency.

As a feasibility test, the public implementation was executed on the
supplied AD dataset using a cut threshold of 403 and automatic
branching-factor selection. It produced an AUC-ROC of
\(0.7742\pm0.0033\) and an average-precision score of
\(0.4460\pm0.0259\). Training and testing took approximately \(19.97\)
and \(38.67\) seconds, respectively. These results are close to the
published AD results. The experiment used Python 3.11.4 and NumPy
1.26.4; NumPy was restricted below version 2.0 because the implementation
uses the removed \texttt{np.mat} function. This successful run provides
initial evidence that the project is technically feasible.

\subsection{Initial LLM-as-Judge Evaluation}

As an exploratory evaluation, an LLM may be asked to judge whether
selected observations appear normal or anomalous from standardised,
human-readable descriptions of their feature values. A small sample will
include correctly detected anomalies, missed anomalies, false alarms and
normal observations. The true labels and OptIForest predictions will be
hidden from the LLM, and its classifications and explanations will be
compared with the ground truth and model output.

This evaluation will supplement rather than replace AUC-ROC, AUC-PR and
error analysis. LLM judgements may be affected by prompt wording,
domain knowledge and model version, so the final protocol will be refined
after the relevant unit material is covered. Prompts, model settings and
responses will be documented, and no private or personally identifiable
information will be provided to an external model.

\section{Original Datasets}

The paper evaluates OptIForest on 20 processed benchmark datasets
\citep{xiang2023optiforest}. Table~\ref{tab:original-data} groups them by
their general data domain. The datasets vary in sample size, feature
dimension and anomaly prevalence, allowing the method to be tested under
substantially different conditions.

\begin{table*}[t]
\centering
\small
\begin{tabular}{p{0.22\textwidth}p{0.68\textwidth}}
\toprule
General domain & Datasets \\
\midrule
Behavioural and content data &
AD, Campaign, Census, Donors and Spam \\
Health-related data &
Arrhythmia and Cardio \\
Network and cybersecurity data &
Backdoor, KDDCup99, HTTP and SMTP \\
Image-derived data &
CelebA and MNIST \\
Scientific and classification benchmarks &
Cover, Ionosphere, Satellite, Shuttle, Vowel, Waveform and Wine \\
\bottomrule
\end{tabular}
\caption{The 20 original benchmark datasets used in the OptIForest
evaluation.}
\label{tab:original-data}
\end{table*}

Several of these resources were originally created for classification
rather than anomaly detection. Their processed anomaly-detection versions
treat selected classes or rare observations as anomalies and the
remaining observations as normal. Consequently, reproducing the results
requires the exact processed version, feature representation and label
conversion used by the authors, rather than merely a dataset with the
same name.

The repository directly supplies only the processed AD dataset. The
remaining datasets must be located through the public sources identified
by the authors, including the UCI Machine Learning
Repository,\footnote{\url{https://archive.ics.uci.edu/}}
Kaggle,\footnote{\url{https://www.kaggle.com/datasets}}
and public anomaly-detection benchmark repositories. For every dataset
used in the final project, its original release or documentation will be
cited and its download location, licence, dimensions, anomaly count and
preprocessing steps will be recorded. If an exact published version
cannot be identified, that dataset will not be presented as an exact
reproduction; it may instead be reported as an extension using a
documented alternative version.

\section{Proposed New Data}

\subsection{Candidate Existing Datasets}
The extension will consider five labelled anomaly-detection datasets that
were not included in the original OptIForest experiments. Labelled data
are preferred because they allow AUC-ROC and AUC-PR to be calculated
consistently with the original evaluation.

\textbf{Credit Card Fraud Detection} contains European credit-card
transactions represented by numerical features. Fraudulent transactions
form a very small minority of the observations, making this dataset
suitable for examining performance under severe class imbalance. It is
publicly available through Kaggle.\footnote{\url{https://www.kaggle.com/datasets/mlg-ulb/creditcardfraud}}

\textbf{Mammography} contains six numerical features derived from
mammogram observations. The minority class represents unusual
observations and can be treated as the anomaly class. Its relatively low
anomaly rate makes it useful for comparing AUC-ROC with AUC-PR.

\textbf{Thyroid} contains numerical medical measurements and labelled
abnormal thyroid cases. It offers a small-to-medium-sized experiment that
should be computationally manageable while still presenting a strongly
imbalanced detection problem.

\textbf{Musk} contains high-dimensional numerical measurements of
molecular conformations. Its large number of features provides an
opportunity to test whether OptIForest remains effective when the feature
dimension is much higher than in the supplied AD dataset.

\textbf{Glass} contains chemical measurements of glass samples. One rare
class is commonly treated as the anomaly class. Although the dataset is
small, it is useful for investigating whether the algorithm's results
become unstable when only a limited number of anomalous observations are
available.

The Mammography, Thyroid, Musk and Glass anomaly-detection versions are
available from established public anomaly-detection benchmark
collections such as ODDS and the Outlier Detection DataSets
repository.\footnote{\url{http://odds.cs.stonybrook.edu/}}
These five datasets were selected to provide variation in sample size,
feature dimension, domain and anomaly rate. Their final inclusion will
depend on confirming their availability, licence conditions, label
definitions and compatibility with the OptIForest input format.

\subsection{Feasibility Assessment}
Table~\ref{tab:candidate-data} provides an initial feasibility comparison.
The dataset statistics will be verified after downloading the final
versions because different public repositories may contain differently
processed copies.

\begin{table*}[t]
\centering
\small
\begin{tabular}{lrrrrl}
\toprule
Dataset & Observations & Features & Anomalies & Approx. rate & Initial assessment \\
\midrule
Credit Card Fraud & 284,807 & 30 & 492 & 0.17\% & Feasible but relatively large \\
Mammography & 11,183 & 6 & 260 & 2.32\% & Highly feasible \\
Thyroid & 3,772 & 6 & 93 & 2.47\% & Highly feasible \\
Musk & 3,062 & 166 & 97 & 3.17\% & Feasible; high dimensionality \\
Glass & 214 & 9 & 9 & 4.21\% & Feasible but statistically unstable \\
\bottomrule
\end{tabular}
\caption{Preliminary feasibility assessment of candidate existing
datasets. Counts are subject to verification against the downloaded
versions.}
\label{tab:candidate-data}
\end{table*}

Mammography and Thyroid are the strongest initial candidates because
they are small enough to support repeated experiments and contain clear
binary anomaly labels. They should require only limited preprocessing,
such as checking missing values, separating labels from features and
standardising the file format expected by the OptIForest implementation.

Musk is also computationally manageable, but its 166 features may
increase training time and memory usage. This makes it useful for testing
high-dimensional performance. Credit Card Fraud is considerably larger
and may require longer execution time. If running the complete dataset is
impractical, a reproducible sampling procedure will be investigated while
preserving all anomalies and documenting the changed anomaly rate.

Glass requires minimal computing resources but contains very few
anomalies. Its evaluation scores may therefore change considerably when
only one observation is ranked differently. It will be treated as a
supplementary stress test rather than the primary extension dataset.

Before final selection, each dataset will be checked for download
availability, usage conditions, feature definitions, missing values,
duplicate observations and label conventions. No personally identifying
information will be retained. All preprocessing will be implemented in
Python scripts and documented in the project repository. Based on this
initial assessment, Mammography and Thyroid will be prioritised, followed
by Musk, subject to successful compatibility testing.

\subsection{Newly Constructed Dataset}
In addition to the existing datasets, a controlled synthetic dataset will
be constructed to test OptIForest under conditions where the source and
type of each anomaly are known. A Python script will generate the data so
that the complete process can be reproduced using a recorded random seed.

The initial dataset will contain approximately 5,000 normal observations
and 100 anomalous observations. Normal observations will be sampled from
multiple Gaussian clusters with different centres, variances and numbers
of observations. This structure is intended to represent a population
containing several legitimate groups rather than one uniform normal
distribution.

Two types of anomalies will be introduced. \textbf{Global anomalies} will
be generated far outside the normal clusters and should therefore be
relatively easy to isolate. \textbf{Local anomalies} will be positioned
near a normal cluster but in regions of low local density. These cases
should be more difficult because their values may appear reasonable when
the overall dataset is considered, but unusual relative to their nearest
cluster.

Every generated row will include a unique identifier, numerical features,
a binary ground-truth label and an anomaly-type label. The anomaly-type
label will be used only for detailed analysis and will not be supplied to
the unsupervised model. The data-generation script will save both the
dataset and a metadata file recording the random seed and generation
parameters.

Several controlled versions will be created by varying dimensionality,
cluster overlap and anomaly rate. For example, experiments may compare
two-dimensional and ten-dimensional data, low and high cluster overlap,
and anomaly rates of approximately \(0.5\%\), \(1\%\) and \(2\%\). These
variations will help determine whether the reported advantage of
OptIForest continues under different data conditions.

The synthetic experiments will compare OptIForest with the standard
Isolation Forest baseline using AUC-ROC, AUC-PR and execution time.
Performance will also be reported separately for global and local
anomalies. A two-dimensional version will be visualised to show the
normal clusters, generated anomalies and model rankings. All scripts,
parameter settings and generated data used in the final analysis will be
included in the project repository.

\section{Replication and Extension Plan}

\subsection{Preliminary Reproduction}

A preliminary reproduction was completed using the supplied AD dataset,
a cut threshold of 403 and automatic branching-factor selection. The
implementation produced mean AUC-ROC and average-precision scores of
\(0.7742\) and \(0.4460\), respectively, which are close to the published
AD results. NumPy was restricted to version 1.26.4 because the code uses
\texttt{np.mat}, which was removed in NumPy 2.0. This successful execution
confirms the project's technical feasibility; complete environment and
timing details are reported in Section~3.5.

\subsection{Planned Experiments}

The project will proceed in four stages. First, a documented Python
environment will be created and selected original experiments will be
reproduced using verified dataset versions and preprocessing procedures.
Second, OptIForest will be compared with standard Isolation Forest using
15 repetitions where computationally practical. AUC-ROC, AUC-PR,
training time and testing time will be recorded.

Third, the effects of branching factor, cut threshold and sampling size
will be investigated systematically. OptIForest will then be evaluated
on selected existing datasets, prioritising Mammography, Thyroid and
Musk, and on the controlled synthetic datasets described in
Section~5.3. Performance on global and local synthetic anomalies will be
examined separately.

Finally, prediction errors and disagreements between OptIForest and the
baseline will be analysed. A small selection may also be used in the
exploratory LLM-as-judge evaluation. Dataset versions, preprocessing
steps, random seeds, parameters and package versions will be recorded.
Raw results will be saved in structured files, and the report's tables
and figures will be generated using reproducible Python scripts.

\subsection{Risks and Limitations}
The main reproducibility risk is that the public repository includes only
the processed AD dataset. Other datasets may be unavailable or may use
versions, preprocessing procedures and anomaly-label definitions that
differ from those used by the authors. Only datasets whose versions can
be identified reliably will be presented as reproductions. Dataset
exclusions and unavoidable preprocessing decisions will be documented.

Additional risks include outdated dependencies, randomised results and
limited computing time. Fixed package versions and recorded random seeds
will be used, and the experimental scope may be reduced using documented
selection criteria if necessary. AUC-PR and error analysis will accompany
AUC-ROC because the latter can appear optimistic for rare anomalies.
Synthetic data may not fully represent real applications, and
LLM-as-judge results may be prompt- or model-dependent. These limitations
will be reported explicitly rather than hidden when interpreting the
results.

\section{Conclusion}
This proposal presents a feasible reproducibility and extension study of
OptIForest for unsupervised anomaly detection. The project will reproduce
selected experiments using the authors' public implementation and
original benchmark data, investigate the effects of important parameters,
and compare OptIForest with the standard Isolation Forest baseline.
Performance will be evaluated using AUC-ROC, AUC-PR, execution time,
repeated experiments and detailed error analysis.

The successful preliminary execution on the supplied AD dataset
demonstrates that the main implementation can be run and can produce
results close to those reported in the paper. The extension will evaluate
the method on additional existing datasets and newly generated synthetic
data containing controlled global and local anomalies. All dependency
versions, preprocessing decisions, parameters, scripts and results will
be documented in the project repository. The study will therefore
determine both whether selected published results are reproducible and
whether the paper's claims continue to hold under new data conditions.

% \bibliographystyle{acl_natbib}
\bibliography{custom}


\end{document}